\documentclass[11pt]{article}

\usepackage[a4paper,margin=1in]{geometry}
\usepackage{amsmath,amssymb}
\usepackage{hyperref}

\hypersetup{
    colorlinks=true,
    linkcolor=blue,
    citecolor=blue,
    urlcolor=blue
}

\title{Lower Semicontinuity}
\author{}
\date{}

\begin{document}

\maketitle

\section*{Definition}

Let $X$ be a topological space and let
\[
f:X\to(-\infty,\infty]
\]
be a function. The function $f$ is lower semicontinuous if, for every
$a\in\mathbb R$, the superlevel set
\[
\{x\in X : f(x)>a\}
\]
is open.

Equivalently, the sublevel sets
\[
\{x\in X : f(x)\le a\}
\]
are closed whenever this formulation is appropriate for the topology under
consideration.

\section*{Sequential Characterization}

If $X$ is a metric space, then $f$ is lower semicontinuous if and only if for
every sequence $x_n\to x$,
\[
f(x)
\le
\liminf_{n\to\infty} f(x_n).
\]

\section*{Interpretation}

Lower semicontinuity means that the function cannot jump downward in the
limit. It may jump upward, but limiting values cannot be strictly smaller than
nearby limiting costs.

\end{document}
